% Algebra 2 Lesson 4-1 Lesson Quiz, short form (from a2_4-1_TE.tex "4-1 Lesson Quiz" block).
% Items 1-3 (the inverse variation items Example 3 teaches) keep the Savvas numbers; items 4-5 (reciprocal-function
% asymptote and translation, Example 5, cut from the fast packet) are left off and listed in the key.
\documentclass[11pt]{article}
\usepackage[margin=0.7in]{geometry}
\usepackage{amsmath,amssymb}
\usepackage{enumitem}
\usepackage{array}
\usepackage{tikz}
\usepackage{pgfplots}
\pgfplotsset{compat=1.18}
\usepackage{fancyhdr}
\usepackage{parskip}
\setlist[enumerate]{leftmargin=*, itemsep=6pt}
\pagestyle{fancy}
\fancyhf{}
\fancyfoot[L]{\small\textit{en}Vision\textsuperscript{\textregistered} Algebra 2}
\fancyfoot[C]{\small 4-1 Lesson Quiz, short form}
\fancyfoot[R]{\small\thepage}
\renewcommand{\headrulewidth}{0pt}
\newcommand{\choice}[1]{\tikz[baseline=-0.6ex]\node[draw, circle, inner sep=0.6pt, minimum size=3.3mm, font=\scriptsize\bfseries] {#1};\;}
\newcommand{\checkbox}{$\square$\;}
\pgfplotsset{quizgrid/.style={
  axis lines=middle, grid=both, grid style={gray!45},
  ticklabel style={font=\scriptsize, fill=white, inner sep=1pt},
  xlabel={$x$}, ylabel={$y$}, axis line style={thick}, clip=true}}

\begin{document}

\noindent\textbf{Name}\; \rule{6cm}{0.4pt} \hfill Period \rule{1.2cm}{0.4pt} \quad Date \rule{2.2cm}{0.4pt}

\vspace{4pt}\noindent\rule{\textwidth}{0.6pt}

\vspace{2pt}\noindent{\Large\textbf{4-1 \; Lesson Quiz: Inverse Variation and the Reciprocal Function}} \hfill {\small short form}

\vspace{4pt}\noindent\rule{\textwidth}{0.6pt}

\vspace{6pt}
\textbf{1.} The chart represents an inverse variation. Find the value of $p$.

\vspace{4pt}
\begin{center}
{\renewcommand{\arraystretch}{1.5}\begin{tabular}{|c|*{4}{>{\centering\arraybackslash}p{1.4cm}|}}
\hline
$x$ & 3 & 10 & 15 & 30 \\ \hline
$y$ & 5 & $p$ & 1 & 0.5 \\ \hline
\end{tabular}}
\end{center}

Show your work. The constant of variation is $k=$ \rule{2cm}{0.4pt} because \rule{6cm}{0.4pt}

\vspace{14pt}\rule{\linewidth}{0.4pt}

\vspace{14pt}\rule{\linewidth}{0.4pt}

\vspace{6pt}
$p=$ \rule{2.5cm}{0.4pt}

\vspace{8pt}\noindent\rule{\textwidth}{0.4pt}

\vspace{6pt}
\textbf{2.} Write an equation for the inverse variation represented by the table.

\vspace{4pt}
\begin{minipage}[t]{0.52\textwidth}
\vspace{0pt}
{\renewcommand{\arraystretch}{1.6}\begin{tabular}{|c|*{4}{>{\centering\arraybackslash}p{1.1cm}|}}
\hline
$x$ & $-3$ & $-1$ & $\frac{1}{2}$ & $\frac{2}{3}$ \\ \hline
$y$ & $4$ & $12$ & $-24$ & $-18$ \\ \hline
\end{tabular}}
\end{minipage}\hfill
\begin{minipage}[t]{0.44\textwidth}
\vspace{0pt}
\begin{itemize}[leftmargin=1.2em, itemsep=6pt]
\item[] \choice{A} $y=\dfrac{x}{-12}$ \qquad\quad \choice{B} $y=\dfrac{-x}{12}$
\item[] \choice{C} $y=\dfrac{12}{x}$ \qquad\quad\;\; \choice{D} $y=\dfrac{-12}{x}$
\end{itemize}
\end{minipage}

\vspace{8pt}
Show one product $xy$ that supports your choice: \rule{6cm}{0.4pt}

\vspace{8pt}\noindent\rule{\textwidth}{0.4pt}

\vspace{6pt}
\textbf{3.} Three students can wash a car in 16 minutes. If the time varies inversely with the number of students washing the car, how many minutes will it take two students to complete that same job?

\vspace{4pt}
\begin{itemize}[leftmargin=1.2em, itemsep=5pt]
\item[] \choice{A} 20 \qquad\qquad \choice{B} 22 \qquad\qquad \choice{C} 24 \qquad\qquad \choice{D} 26
\end{itemize}

\vspace{4pt}
Show your work. Equation: $t=$ \rule{3cm}{0.4pt} \qquad $k=$ \rule{2.5cm}{0.4pt}

\vspace{14pt}\rule{\linewidth}{0.4pt}

\vspace{14pt}\rule{\linewidth}{0.4pt}

\vfill
\noindent\textbf{Scoring:} Question 1 and Question 3 are 10 points each. Question 2 is a bonus of up to 5 points. Total out of 20.

\newpage
\noindent{\LARGE\textbf{4-1 Lesson Quiz, short form: answer key}}

\vspace{4pt}\noindent\rule{\textwidth}{0.6pt}

\begin{enumerate}
\item[\textbf{1.}] \textbf{$p=1.5$.} $k=xy=3(5)=15$ (also $15(1)=15$ and $30(0.5)=15$), so $10p=15$ and $p=\frac{15}{10}=1.5$. Savvas quiz item 1, DOK 1, Realize skill F74. Taught by Example 3 (find $k$ from a known pair, then use it) and rehearsed by the packet's Practice table. \textbf{10 points:} 4 for $k=15$ with a product shown, 4 for $p=1.5$, 2 for the ``because'' (the product $xy$ is the same for every pair).
\item[\textbf{2.}] \textbf{D, $y=\dfrac{-12}{x}$.} Every product is $-12$: $(-3)(4)=(-1)(12)=\frac12(-24)=\frac23(-18)=-12$. Savvas quiz item 2, DOK 2, Realize skill F74. Taught by Example 3 Step 1 ($k$ from one pair) and the packet's Practice table (negative $k$). Distractors A and B are direct variations ($y$ a multiple of $x$); C has the wrong sign. \textbf{Bonus 5:} 5 for D with one product shown; 3 for D alone.
\item[\textbf{3.}] \textbf{C, 24 minutes.} With $n$ students, $t=\frac{k}{n}$; $16=\frac{k}{3}$, so $k=48$; $t=\frac{48}{2}=24$. Savvas quiz item 3, DOK 1, Realize skill F74. Taught by Example 3 (the same two steps in a context) and Try It 3. \textbf{10 points:} 4 for the equation $t=\frac{48}{n}$ with $k=48$, 4 for C, 2 for the work shown ($t=\frac{48}{2}$).
\end{enumerate}

\textbf{Items left off this form} (available for a retake or Form B; both are Example 5, which the fast packet cuts): Savvas item 4, the horizontal asymptote of $y=\frac{1}{x+5}-4$ (answer $y=-4$, DOK 2, F74); Savvas item 5, $y=\frac1x$ translated 3 units up and 2 units to the left (answer A, $y=\frac{1}{x+2}+3$, DOK 2, F74). The TE's Lesson Quiz has no graphs, so nothing was redrawn.

\textbf{IXL replacement:} SmartScore 70 or above on PNY and PXE replaces this quiz with 20/20.

\end{document}
